\unnumberedchapter{ЗАКЛЮЧЕНИЕ}

В результате дипломного проекта разработано веб-приложение ChessView для шахматного онлайн-взаимодействия пользователей. Приложение включает авторизацию, пользовательские профили, live-партии, серверную проверку ходов, WebSocket-события, чат, историю партий, анализ позиций, шахматные задачи, турниры, запланированные матчи, демонстрационный платежный сценарий и отдельную административную панель.

В основной части рассмотрена проблематика разработки шахматного веб-приложения. Были выделены ключевые сложности: синхронизация состояния партии, серверная проверка ходов, учет времени, восстановление соединения, хранение истории, разделение REST и WebSocket, разграничение пользовательских и административных функций. Также выполнено сравнение аналогов, описаны инструменты разработки и определены требования к серверу для запуска приложения.

В технической части подробно описана структура программного обеспечения. Backend разделен на домены и слои presentation, application, domain и infrastructure. Пользовательский frontend организован по слоям app, pages, widgets, features, entities и shared. Административный frontend вынесен в отдельное приложение и работает через admin API.

Отдельно описана база данных. Для нее построена ERD-модель, разбитая на несколько небольших схем, и приведено описание всех таблиц и всех полей. Такой формат делает структуру БД понятной: можно увидеть как общие связи между сущностями, так и конкретное назначение каждой колонки.

API описано отдельно от клиентской функциональности. REST API используется для управляемых операций, а WebSocket API используется для событий live-партии, очереди, чата и RTC signaling. Функциональность клиента и административной панели описана как инструкция пользователя и администратора, чтобы было понятно не только из чего состоит проект, но и как им пользоваться.

Поставленная цель достигнута. ChessView является воспроизводимым full-stack проектом, который можно запустить локально через Docker Compose или раздельно в режиме разработки. Проект демонстрирует практическое применение современных инструментов веб-разработки и может быть расширен за счет распределенного WebSocket fanout, промышленного хранения media-файлов, внешнего платежного провайдера и дополнительных античит-механизмов.
